  \chapter{Thermal Effects on Breakdown and Transition}\label{chapter:breakdown}

Turbulence development in hypersonic boundary layers proceeds through a sequence beginning with linear instability growth and culminating in nonlinear breakdown. Whilst \gls{LST} captures the second Mack mode during early amplification, the subsequent nonlinear dynamics remain difficult to predict with comparable accuracy. Breakdown of coherent structures into turbulence determines the heating, frictional drag, and aerothermal performance of hypersonic vehicles. A characteristic feature is the overshoot in surface quantities such as skin-friction coefficient and Stanton number. Rather than the gradual increase predicted by linear theory, both exceed their fully turbulent values during transition before relaxing downstream \citep{Holden1972}. This behaviour arises from abrupt, localised enhancement of mixing and momentum transport during vortex breakup, which injects turbulent-like fluctuations into the near-wall region faster than the mean flow equilibrates. This transient amplification of aerodynamic heating poses a significant challenge for \gls{TPS} design and risks localised material failure. Understanding this response remains critical for hypersonic vehicle design across a wide range of flight regimes.

Second-mode instabilities are the dominant mechanism in high-Mach-number boundary layers. Controlled-forcing experiments using harmonic point sources confirmed the second mode as the primary instability and identified subharmonic resonance as a breakdown route \citep{Shiplyuk2003,Bountin2008}. More recent work in the Purdue quiet tunnel identified fundamental resonance as the dominant transition pathway on a Mach~6 flared cone \citep{Chynoweth2014} and established amplitude thresholds for turbulent spot formation \citep{Casper2011,Casper2012}. Low-enthalpy studies, including Mach~21 nitrogen-tunnel experiments \citep{Mironov2000}, revealed four-wave parametric resonance between fundamental and harmonic modes. However, such experiments remain limited by quiet-tunnel availability, perfect-gas conditions, and incomplete access to nonlinear dynamics. High-fidelity \gls{DNS} is necessary to systematically assess thermochemical nonequilibrium effects and wall thermal variations.

Temporal and spatial \gls{DNS} studies have investigated fundamental transition mechanisms in compressible hypersonic boundary layers. For instance, \citet{Mayer2011} computed Mach~3 transition triggered by oblique breakdown, a pathway first identified by \citet{Fasel1993}. Such simulations typically employ either broadband forcing across multiple frequencies or selective excitation of dominant energy modes \citep{AlamSandham2000,Bhaskaran2011}. Whilst the latter strategy does not replicate realistic flight or noisy wind tunnels, it efficiently isolates underlying physics. Fundamental and subharmonic resonance originate from amplification of a two-dimensional instability wave. Once this wave reaches finite amplitude, secondary oblique instabilities emerge, growing either at the same frequency (fundamental resonance) or at half the frequency (subharmonic resonance). The current study examines the transition mechanism through the second excitation method, building on \Cref{chapter:disturbance_growth}. Fundamental \gls{DNS} studies of turbulent boundary-layer behaviour in hypersonic flows were initiated by \citet{Martin2007} and advanced by \citet{Duan2011}, who examined the influence of wall temperature and Mach number on mean flow statistics in cold hypersonic flows. \citet{Franko2013} performed a spatial \gls{DNS} of a Mach~6 \gls{ZPG} boundary layer, examining the skin-friction and heat-transfer overshoot observed in experiments and identifying first-mode oblique breakdown as a plausible transition path. An extensive database of cooled boundary layers up to Mach~14 was later generated by \citet{Zhangetal2018}.

Fewer studies have incorporated real gas effects. \citet{Stemmer2005} simulated hypersonic boundary layer breakdown for reactive and non-reactive flow in a Mach~20 boundary layer. \citet{Marxenetal2013,Marxenetal2014} incorporated fundamental resonance and nonlinear modes in a Mach~10 adiabatic flat-plate flow, concluding that chemical reactions had no significant effect on growth rates. \citet{DiRenzo2021} performed a complete transition study incorporating chemical nonequilibrium and observed good qualitative agreement with previous works. \citet{Passiatoreetal2021} provided a comprehensive analysis of breakdown to turbulence under subharmonic second-mode excitation, identifying that fundamental mode saturation and energy transfer to oblique subharmonic modes lead to $\Lambda$-vortex formation. Whilst finite-rate chemistry had little effect on initial transition stages, the endothermic nature of later stages drained modal energy and delayed transition by approximately 7\%.

Several \gls{DNS} studies have examined real-gas effects on hypersonic boundary layers, though most focused on chemical nonequilibrium. \citet{DiRenzo2021} investigated transition in a Mach~10 hypersonic boundary layer at suborbital enthalpy, capturing the sequence of second-mode amplification, resonance with oblique harmonics, and breakdown, yielding four- to five-fold increases in wall friction and heating. That work also considered finite-rate chemistry in cooled turbulent boundary layers, whilst \citet{Passiatoreetal2021} extended the framework to pseudo-adiabatic walls to isolate chemical nonequilibrium effects. A significant limitation was the neglect of thermal nonequilibrium. More recently, \citet{Williamsetal2025} examined chemistry--turbulence coupling and showed that turbulence drives variations in chemical source terms whilst having minimal impact on turbulence structure. Most studies have been limited to fully thermochemical equilibrium conditions at unrealistically high Mach numbers, leaving understanding of isolated thermal effects incomplete.

This chapter extends the \gls{LST} analysis to the breakdown stage across \gls{CPG}, \gls{TEQ}, and \gls{TNE} formulations. All cases are treated as chemically frozen, as the selected conditions do not trigger significant chemical activity \citep{BitterShepherd2015}. At high-enthalpy conditions, deviations from equilibrium become important, and it is necessary to assess how different thermal models affect breakdown onset and progression. Attention is given to the interplay between linear amplification, nonlinear resonance, secondary instability growth, and turbulent spot formation. Emphasis is placed on the overshoot behaviour of skin friction and heat transfer, the influence of cold-wall conditions in accelerating breakdown, and the collapse of Mack second modes into fully developed turbulence across the thermal sates. The analysis further examines the wall-temperature response of \gls{TNE} flows, where vibrational nonequilibrium modifies breakdown dynamics compared to \gls{CPG} and \gls{TEQ} cases. These results provide a detailed characterisation of the nonlinear transition regime under different thermal models, Mach numbers and wall conditions, highlighting the consequences for aerothermodynamic design and the challenges in developing predictive tools.

\section{Numerical setup and flow definitions}\label{section:breakdown:numericalsetup}


This study examines spatially developing \gls{ZPG} flat-plate boundary layers formed behind a shock wave of a wedge placed in a Mach 10 freestream at 36 km altitude, conditions representative of the heat-rate limit along a hypersonic glide-vehicle trajectory~\citep{Rizvi2015} as detailed in \Cref{section:disturbances:computationdomain}. Post-shock conditions of Mach 6 and 8 are investigated. The edge conditions are $p_e = 4.91~\mathrm{kPa}$ and $T_e = 620~\mathrm{K}$ for Mach 6, and $p_e = 1.78~\mathrm{kPa}$ and $T_e = 365~\mathrm{K}$ for Mach 8 as mentioned in~\Cref{table:flight_conditions}. The governing equations are the compressible Navier--Stokes equations under thermal nonequilibrium and chemically frozen assumptions. The computational domain is a rectangular box with streamwise extents of $800\delta^*_0$ (Mach 6) and $1600\delta^*_0$ (Mach 8), wall-normal height of $100\delta^*_0$, and spanwise width of $20\delta^*_0$. The grid is uniform in the streamwise and spanwise directions, with wall-normal spacing generated via hyperbolic tangent stretching, $y = L_y \sinh(s_1 \xi)/\sinh(s_1)$ for uniformly distributed $\xi \in [0,1]$ and stretching factor $s_1 = 5.0$. Grid resolutions are presented in \Cref{tab:domain_grid_dns} and \Cref{tab:domain_grid_iles}.

\begin{table}[hbt!]
\centering
\caption{Grid resolution and thermal conditions for DNS simulations.}
\begin{tabular}{lccccc}
\toprule
Case & Gas Model & $M_e$ & $T_w/T_e$ &  ($N_x \times N_y \times N_z$) & Total points \\
\midrule
\multirow{3}{*}{M6Tw18} & CPG & 6 & 2.90 & $5200 \times 600 \times 275$ & $8.58\times 10^8$ \\
 & TEQ & 6 & 2.90 & $5200 \times 600 \times 275$ & $8.58\times 10^8$ \\
 & TNE & 6 & 2.90 & $5200 \times 600 \times 275$ & $8.58\times 10^8$ \\
\bottomrule
\end{tabular}
\label{tab:domain_grid_dns}
\end{table}
\begin{table}[hbt!]
\centering
\caption{Grid resolution and thermal conditions for ILES simulations.}
\begin{tabular}{lccccc}
\toprule
Case & Gas Model & $M_e$ & $T_w/T_e$ &  ($N_x \times N_y \times N_z$) & Total points \\
\midrule
M6Tw12 & TNE & 6 & 1.93 & $2600 \times 600 \times 205$ & $3.20\times 10^8$ \\
M6Tw18 & TNE & 6 & 2.90 & $2300 \times 500 \times 165$ & $1.84\times 10^8$ \\
M8Tw12 & TNE & 8 & 3.29 & $4600 \times 500 \times 165$ & $3.68\times 10^8$ \\
M8Tw18 & TNE & 8 & 4.93 & $4600 \times 500 \times 165$ & $3.68\times 10^8$ \\
\bottomrule
\end{tabular}
\label{tab:domain_grid_iles}
\end{table}

Inflow boundary-layer profiles are obtained from a locally self-similar framework solved using Chebyshev collocation, ensuring consistent specification across \gls{CPG}, \gls{TEQ}, and \gls{TNE} conditions. For \gls{TNE}, a thermally frozen \gls{TFG} similarity solution provides the inflow as detailed in \Cref{section:similaritysolution}. Conservative variables are initialised and imposed via inlet pressure extrapolation. Outflow and far-field boundaries employ simple extrapolation, whilst spanwise boundaries are periodic. The wall is non-catalytic and isothermal at $T_w = 1200~\text{K}$ (cold) or $T_w = 1800~\text{K}$ (hot). Cold walls enhance second-mode amplification by thinning the boundary layer and destabilising the acoustic spectrum, whilst hot walls stabilise acoustic-trapped disturbances, delaying transition.

The inflow Reynolds number is $Re_{\delta^*_0} = 10{,}000$. The governing equations are discretised using a 5th-order WENO-Z scheme for convective terms and 4th-order central differences for viscous and metric terms, with 3rd-order SSP Runge--Kutta time integration. Transport properties follow the Musawi--Sandham viscosity model~\citep{Musawi2025} and vibrational relaxation uses the Millikan--White formulation with Park constants~\citep{Park1993}. Non-dimensional timesteps are $\Delta t^{*} = 0.012$ (Mach 6) and $\Delta t^{*} = 0.016$ (Mach 8) with a \gls{CFL} of $0.05$ across all cases, ensuring stability whilst resolving high-frequency acoustic waves characteristic of second-mode instability.

\subsection{Disturbance seeding and excitation method}

Transition to turbulence is induced via a suction--blowing forcing strip following \citep{AlamSandham2000}, extended here to a two-dimensional domain with large-amplitude perturbations. A wall mass-flux perturbation is imposed at $x_f$:
\begin{equation}
\rho v_w(x,z,t) \;=\; \rho a_0 \, g(x) \, h(z) \, \sum_{i=1}^{n} \sin(2\pi f_i t),
\end{equation}
where the streamwise envelope is
\begin{equation}
g(x) = \exp\left(-\frac{(x-x_f)^2}{2(\delta^*_0)^2}\right),
\end{equation}
and the spanwise modulation is
\begin{equation}
h(z) = 1 + 0.06 \left[\sin\!\left(\frac{4\pi z}{L_z}\right) + \sin\!\left(\frac{2\pi z}{L_z}\right)\right].
\end{equation}
The Gaussian envelope has width $2\sigma^2 = 4(\delta^*_0)^2$ and is centred at $x_f = 25\delta^*_0$ from the domain origin. The forcing amplitude is $a_0 = 0.1 \, u_e$, sufficiently large to trigger transition within the domain whilst remaining representative of realistic disturbance levels. The frequency set is case-dependent: for Mach 6, $f^{*}_i = \{0.15,\,0.16\}$, and for Mach 8, $f^{*}_i = \{0.10,\,0.11\}$. These frequencies are selected using linear stability theory to ensure second-mode growth reaches its maximum between $x/\delta^*_0 \approx 220$ and $400$ across all cases, coinciding with the forcing location at $x_f \approx 25\delta^*_0$. This alignment excites the most naturally amplified disturbances and sustains their growth within the computational domain.


\begin{figure}[t]
    \centering
    \includegraphics[width=9cm]{resources/figures/breakdown/insert_shapefunction_natural.pdf}
    \caption{Spatial structure of the wall mass-flux forcing function.}
    \label{figure:forcing_shapefunction}
\end{figure}

The spatial structure of the forcing function is shown in \Cref{figure:forcing_shapefunction} with a comparison to~\citep{Franko2013}. Rather than imposing oblique modes, three-dimensionality is introduced through spanwise modulation, which imposes long- and short-wavelength subharmonics ($\lambda_z = L_z/2$ and $L_z$). This approach is justified by the stability analysis in \Cref{figure:stability_plot}, which presents a Mach 6 case under different thermal excitations with adiabatic wall results in the left column and ``hot" wall results in the right column. While first mode instability is weakly presetned for adibatic case in \Cref{figure:stability_plot}, the analysis demonstrates the absence of first-mode instability across all thermal models at the presented wall condition; consequently, oblique modes arising from first-mode interactions need not be artificially introduced. Instead, $\Lambda$-shaped vortices and streaks emerge naturally, isolating the evolution of two-dimensional second-mode disturbances and their secondary instabilities without prescribing an artificial oblique structure.

\begin{figure}[t]
    \centering
    \includegraphics[width=12cm]{resources/figures/breakdown/result_mach6_stability_diagram.pdf}
    \caption{Comparison of Mach 6 stability diagrams at $\delta_{*}^{0}=255$ for (a) CPG, (b) TEQ, and (c) TNE models. For each model, the left column shows the adiabatic-wall case, while the right column corresponds to the isothermal 1800\,K wall condition.}
    \label{figure:stability_plot}
\end{figure}

Each simulation spans four consecutive flow-through times ($t_f = L_x/u_e$). The first two flow-throughs remove initial transients; the solution is then restarted for two further flow-throughs whilst accumulating Favre-averaged statistics. This c aptures the complete sequence from controlled second-mode excitation through nonlinear interactions to turbulence onset.

\section{Influence of thermal excitations on flow behaviour}

The influence of thermal modelling on hypersonic boundary layer behaviour is investigated by comparing \gls{CPG}, \gls{TEQ}, and \gls{TNE} formulations for the Mach 6 hot-wall configuration. Linear stability analyses reveal that thermal modelling substantially affects second-mode growth rates. \gls{TEQ} exhibits the highest growth rates, exceeding those from \gls{TNE} and \gls{CPG}, due to stronger coupling between vibrationally equilibrated energy modes and acoustic disturbances. The \gls{TNE} formulation introduces vibrational relaxation effects that reduce modal amplification, whilst \gls{CPG} yields the weakest instability growth due to the absence of additional energy channels. These differences arise primarily from variations in the base flow. The resulting \gls{LST} predictions provide context for interpreting the nonlinear transition dynamics in the simulations.

Transition is characterised through skin-friction coefficient, Stanton number, mean velocity profiles, and velocity fluctuation components, with onset identified from the sharp rise in skin friction and Stanton number. Reynolds stress components and wall-normal heat flux are examined across the transitional and turbulent regions to assess how the three thermal models differ in turbulence production and identify the mechanisms responsible for heat-transfer overshoot.


\subsection{Early disturbance amplification}
The streamwise evolution of wall-bounded quantities provides diagnostics for characterising the onset and progression of boundary layer transition. The skin-friction coefficient and Stanton number are defined respectively as
\begin{equation}
    C_f = \cfrac{2\tau_w}{\rho_e u_e^2},
\end{equation}
\begin{equation}
    St = \cfrac{q_w}{\rho_e u_e c_{p,tr} (T_{aw}-T_w)},
\end{equation}
where $T_{aw}$ is the adiabatic wall temperature derived from the respective self-similar solution. For the Stanton number, the translational-rotational temperature $T$ is used for \gls{CPG} and \gls{TNE}, whilst \gls{TEQ} employs a $c_{p,trv}$ because of the equilibrium state. The skin-friction coefficient reflects near-wall shear stress and disturbance development, whilst the Stanton number quantifies the nondimensional wall heat flux and thermal boundary-layer response. Both parameters characteristically exhibit overshoot relative to the fully turbulent asymptote during nonlinear breakdown.

For laminar boundary layers, the skin-friction and Stanton number coefficients are derived from self-similarity solutions at each streamwise location to provide a baseline for assessing disturbance growth. Following \citet{White2005}, these are expressed as
\begin{equation}
c_{f,lam} = \sqrt{2} f''(0) \,Re_x^{-1/2}\cfrac{\mu_w \rho_w}{\mu_e \rho_e},
\end{equation}
where $f$ is the dimensionless velocity similarity variables. Both depend on the thermal model; formulations are derived from their respective similarity solutions for \gls{CPG}, \gls{TEQ}, and \gls{TNE}.

For turbulent boundary layers, predictions are benchmarked against the compressible Van Driest II correlation \citep{VanDriest1952,Sandhametal2014}:
\begin{equation}
    c_{f,turb} = \frac{1}{\lambda} \left[ \frac{0.242 \left( \arcsin\alpha + \arcsin\beta \right)}{\log_{10}(c_{f,turb}  Re_x) - 0.76 \log_{10}(T_w/T_e) + 0.41} \right],
\end{equation}
where $\lambda = (\gamma - 1) M_e^2 r/2$, $\alpha = (2a^2 - b)/Q$, $\beta = b/Q$, $Q = \sqrt{b^2 + 4a^2}$, $a = \sqrt{\lambda T_e/T_w}$, and $b = (1+\lambda)T_e/T_w - 1$. The Stanton number for both laminar and turbulent counterparts is obtained via the Reynolds analogy:
\begin{equation}
St = \frac{c_f}{2Pr_{tr}^{2/3}}.
\end{equation}

\begin{figure}[t]
    \centering
    \fontsize{9}{10}\selectfont
    \includegraphics[width=0.65\textwidth]{resources/figures/breakdown/result_wallquantities_cf_st.pdf}
    \caption[Streamwise evolution of spatially temporally averaged $C_f$ and $St$ for the Mach~6, $T_w=1800$ K case. Results are shown for CPG, TEQ, and TNE formulations, alongside laminar reference solutions and the Van Driest~II turbulent correlation.]%
    {Streamwise evolution of spatially temporally averaged $C_f$ and $St$ for the Mach~6, $T_w=1800$ K case. Results are shown for \gls{CPG}, \gls{TEQ}, and \gls{TNE} formulations, alongside laminar reference solutions and the Van Driest~II turbulent correlation.}
    \label{figure:thermal_cfst}
\end{figure}

In the present work, laminar and turbulent reference profiles are shown only for the \gls{TNE} case, to avoid excessive clutter in the plots. Comparisons indicate that the laminar and turbulent trends for \gls{CPG} and \gls{TEQ} are very similar to those for \gls{TNE}, although the \gls{TEQ} profiles lie slightly lower. Therefore, the use of \gls{TNE} profiles as representative references is sufficient.  The skin friction and Stanton number across all cases are shown in \Cref{figure:thermal_cfst}. In the laminar region, all three models closely follow the laminar reference with minor thermal-model differences. A pronounced second-mode resonance peak emerges, with \gls{TEQ} exhibiting a sharp, well-defined peak, \gls{CPG} showing a similar structure delayed by a few displacement thicknesses with slightly higher magnitude, and \gls{TNE} displaying the most attenuated peak. This progressive delay and attenuation reflects different base flows and vibrational relaxation extracting energy from the acoustic modes driving second-mode growth, consistent with~\citep{DiRenzo2021} though not observed in~\cite{Passiatore2021} or \cite{Williamsetal2025}.

Nonlinear breakdown then accelerates, with both quantities overshooting the turbulent asymptote. \gls{TEQ} exhibits the most dramatic overshoot followed by \gls{CPG} further downstream, whilst \gls{TNE} shows considerably subdued overshoot with a more gradual transition to turbulence. The transition hierarchy is clear: \gls{TEQ} transitions earliest, \gls{CPG} transitions downstream, and \gls{TNE} transitions latest, delayed by approximately $150\delta^*$ relative to \gls{TEQ}. The contrasting dynamics reflect both the base-flow effects and the thermal model differences; \gls{TEQ} and \gls{CPG} exhibit sharp transitions with pronounced overshoots followed by rapid relaxation, whilst \gls{TNE} shows a prolonged, gentler approach with suppressed overshoot, demonstrating the stabilising influence of thermal nonequilibrium. Selected flow quantities and grid parameters at representative stations throughout the transition process are presented in \Cref{table:global_quantities_models}, documenting streamwise evolution across five characteristic phases: near resonance of the two-dimensional second-mode, secondary instability growth, imminent breakdown, end of transition, and the fully turbulent state. Boundary-layer integral parameters including the shape factor $H$ and momentum-thickness Reynolds number $Re_\theta$ characterise velocity profile development, showing progressive steepening as the flow transitions from laminar to turbulent. Thermal quantities including maximum temperature rise $T_{max}/T_e$ quantify thermal excitation effects. In \gls{TNE} flows, the energy balances $B_q = q_w/(\rho_w u_\tau h_w)$ and $B_{q,\mathrm{v}} = q_{w,\mathrm{v}}/(\rho_w u_\tau h_w)$ where $h_w$ is the enthalpy at the wall. These parameters reveal opposing signs characteristic of energy partitioning between the modes. A clear delay in transition stations is evident across the thermochemical models, with \gls{TNE} consistently requiring greater streamwise distances to progress through successive phases compared to the perfect-gas assumption. The friction Reynolds number $Re_\tau$ and friction Mach number $M_\tau$ increase substantially during breakdown, reflecting intensified near-wall velocity gradients as coherent structures develop and break down. Grid spacings satisfy established \gls{DNS} resolution requirements for capturing near-wall viscous structures and streamwise coherent features, consistent with previous studies~\citep{PoggieEtAl2015}. Similar resolutions were recently achieved using the TENO scheme~\citep{Williamsetal2025}.

\newpage
\begin{sidewaystable}[H]
\centering
% \small
\caption{Streamwise evolution of span-averaged and time-averaged dimensionless flow quantities across the transition process for Mach 6, $T_w = 1800$ K CPG, TEQ and TNE simulations.}
\begin{tabular}{lccc ccc ccc ccc ccc}
\toprule
 & \multicolumn{3}{c}{Near resonance of 2D mode} 
 & \multicolumn{3}{c}{Secondary instability} 
 & \multicolumn{3}{c}{Imminent breakdown} 
 & \multicolumn{3}{c}{End of transition} 
 & \multicolumn{3}{c}{Turbulent} \\
\cmidrule(lr){2-4} \cmidrule(lr){5-7} \cmidrule(lr){8-10} \cmidrule(lr){11-13} \cmidrule(lr){14-16}
 & CPG & TEQ & TNE & CPG & TEQ & TNE & CPG & TEQ & TNE & CPG & TEQ & TNE & CPG & TEQ & TNE \\
\midrule
$x/\delta^*_0$ & 256 & 252 & 292 & 312 & 291 & 361 & 400 & 337 & 436 & 505 & 388 & 505 & 789 & 600 & 600 \\
$Re_x \times 10^{-6}$ & 4.16 & 4.29 & 4.53 & 4.72 & 4.68 & 5.22 & 5.60 & 5.14 & 5.97 & 6.65 & 5.65 & 6.66 & 9.49 & 7.77 & 7.61 \\
$Re_\tau$ & 138 & 145 & 189 & 160 & 151 & 166 & 387 & 175 & 240 & 514 & 254 & 514 & 913 & 687 & 824 \\
$Re^*_\delta \times 10^{-4}$ & 1.61 & 1.56 & 1.68 & 1.72 & 1.63 & 1.80 & 1.87 & 1.71 & 1.93 & 2.04 & 1.79 & 2.04 & 2.43 & 2.10 & 2.18 \\
$Re_\theta$ & 1177 & 1132 & 1211 & 1365 & 1200 & 1333 & 1815 & 1294 & 1537 & 2469 & 1437 & 1809 & 4974 & 2123 & 2323 \\
$H$ & 13.62 & 13.22 & 13.96 & 12.20 & 13.11 & 14.00 & 10.64 & 12.59 & 12.69 & 10.33 & 11.73 & 11.03 & 10.24 & 10.04 & 10.71 \\
$M_{\tau,\mathrm{max}}$ & 0.10 & 0.08 & 0.09 & 0.11 & 0.08 & 0.07 & 0.16 & 0.08 & 0.09 & 0.15 & 0.09 & 0.14 & 0.14 & 0.11 & 0.15 \\
$B_{q,\mathrm{tr}} \times -10^{-2}$ & 7.17 & 2.93 & 5.11 & 6.23 & 2.88 & 3.78 & 8.98 & 3.00 & 4.72 & 8.24 & 3.41 & 6.99 & 7.37 & 4.21 & 7.23 \\
$B_{q,\mathrm{v}} \times 10^{-3}$ & 0.00 & $-5.87$ & 5.58 & 0.00 & $-5.77$ & 5.79 & 0.00 & $-6.01$ & 5.64 & 0.00 & $-6.83$ & 5.69 & 0.00 & $-8.45$ & 5.43 \\
$T_{max}/T_e$ & 3.70 & 3.36 & 3.73 & 3.70 & 3.37 & 3.67 & 3.69 & 3.40 & 3.75 & 3.59 & 3.48 & 3.81 & 3.56 & 3.51 & 3.72 \\
$\delta_{M_\tau}/\delta^{*}_{0}$ & 0.88 & 0.90 & 0.90 & 0.91 & 0.90 & 0.90 & 0.47 & 0.89 & 0.86 & 0.37 & 0.87 & 0.67 & 0.48 & 0.66 & 0.33 \\
\midrule
$\Delta x^+$ & 7.72 & 5.49 & 6.56 & 7.99 & 5.47 & 5.31 & 12.14 & 5.88 & 6.72 & 11.47 & 6.67 & 10.05 & 10.66 & 8.35 & 10.97 \\
$\Delta y_w^+$ & 0.68 & 0.48 & 0.58 & 0.70 & 0.48 & 0.47 & 1.07 & 0.52 & 0.59 & 1.01 & 0.59 & 0.88 & 0.94 & 0.73 & 0.96 \\
$\Delta y^+_{99}$ & 1.54 & 1.52 & 1.97 & 1.74 & 1.58 & 1.72 & 4.00 & 1.81 & 2.46 & 5.22 & 2.59 & 5.20 & 9.15 & 6.89 & 8.27 \\
$\Delta z^+$ & 3.59 & 2.55 & 3.05 & 3.71 & 2.54 & 2.47 & 5.63 & 2.73 & 3.12 & 5.33 & 3.10 & 4.67 & 4.95 & 3.88 & 5.09 \\
$L_z^+$ & 1004 & 714 & 853 & 1038 & 711 & 690 & 1578 & 764 & 873 & 1491 & 867 & 1306 & 1385 & 1085 & 1426 \\
\bottomrule
\end{tabular}
\label{table:global_quantities_models}
\end{sidewaystable}
\newpage

\begin{figure}[H]
    \centering
    \includegraphics[width=22.3cm,angle=-90]{resources/figures/breakdown/result_dns_thermal_stanton_contour.pdf}
    \caption{Time-averaged wall Stanton number ($St$) during boundary layer transition for Mach 6, $T_w = 1800$ K case (a) CPG, (b) TEQ, and (c) TNE models.}
    \label{figure:thermal_spanwise_st}
\end{figure}

The spanwise organisation of heat transfer during transition, as shown in \Cref{figure:thermal_spanwise_st}, reveals noticeable differences across the three thermochemical models. The figure displays contours of Stanton number ($St$) as a function of streamwise position and spanwise coordinate, spanning the transitional region between $x/\delta^*_0 = 200$ and $550$. Streak-like structures emerge as transition progresses, redistributing heat flux into alternating high- and low-intensity bands that mark secondary instability growth and rapid breakdown, consistent with observations in prior transition studies \citep{Franko2013}. The spatial organisation and intensity of these streaks differ markedly across models. \gls{TEQ} develops well-defined streaks early in transition, whilst \gls{CPG} exhibits comparable streak organisation with marginally higher peak intensities during breakdown. Both models maintain organised patterns well into the turbulent region. \gls{TNE} exhibits markedly delayed and attenuated streak development, with high-temperature regions emerging considerably further downstream and remaining less intense than in either \gls{CPG} or \gls{TEQ}. By the end of transition, streak organisation weakens across all models and wall heat flux becomes more uniformly distributed, consistent with fully developed turbulence. This progression from well-organised streaks in \gls{TEQ} and \gls{CPG} to attenuated, delayed streaks in \gls{TNE} reflects the stabilising influence of vibrational nonequilibrium on coherent-structure formation and wall heat transfer.


\subsection{Transition structures}

The initial phase of transition is characterised by amplification of two-dimensional disturbances dominated by the Mack second mode. Linear stability theory predictions presented in \Cref{chapter:disturbance_growth} demonstrate that the second mode exhibits the largest amplification rates at the frequencies present in this flow. \Cref{figure:lstdns_amplitude} presents linear amplitude growth curves for the representative cases studied herein. The forced frequencies $f = 0.15$ and $f = 0.16$ exhibit peak amplitude growth at distinct streamwise locations, with the observed resonance peaks in the wall-bounded quantities aligning closely with these linear growth regions. This alignment confirms the transition from linear to weakly nonlinear dynamics as energy from the primary modes redistributes into secondary instabilities.

Instantaneous snapshots reveal the coherent structures driving transition. Q-criterion isosurfaces coloured by total density $\rho / \rho_e$, with pressure contours $p/(\rho_e u_e^2)$ in the background (\Cref{figure:transition_isosurfaces}), show how instability modes develop and interact across the three thermal models. Initially, two-dimensional Mack modes develop as elongated streamwise-aligned structures. These span-aligned structures undergo resonance amplification at the linear growth regions predicted by linear stability theory, transferring energy from the primary mode into secondary instabilities. The secondary structures drive streak formation~\citep{Klebanoff1962} and localised heat-transfer intensification seen in the spanwise distributions, with the strength and organisation of these streaks determining the character of the heat-flux patterns during breakdown.

\begin{figure}[H]
    \centering
    \includegraphics[width=12.6cm,angle=0]{resources/figures/breakdown/result_dns_isosurfaces.pdf}
    \caption{Q-criterion isosurfaces coloured by density during transition for (a) CPG, (b) TEQ, and (c) TNE models. Pressure field shown in background plane. Isosurfaces reveal two-dimensional Mack mode development in the initial domain region and subsequent three-dimensional modal interactions driving breakdown.}
    \label{figure:transition_isosurfaces}
\end{figure}

The three-dimensional boundary-layer structure during transition is illustrated through instantaneous streamwise Q-criterion isosurfaces, which reveal the development of low-speed and high-speed streaks characteristic of secondary instability growth across the three thermal models. H-type transition, the most commonly observed phenomenon in hypersonic boundary layers~\citep{DeTullio2013}, is driven by destabilisation of two-dimensional Mack modes~\citep{Fedorov2011} and their interaction with oblique second modes~\citep{Franko2013}. The \gls{TEQ} case exhibits pronounced, coherent streak formation beginning at comparatively early streamwise locations, with alternating regions of reduced and elevated streamwise velocity apparent at early downstream positions. This vigorous streak development reflects rapid secondary instability growth driven by sustained forcing. The \gls{CPG} case shows similar streak structures but emerging further downstream, reflecting slightly delayed transition onset. In contrast, the \gls{TNE} case exhibits considerably delayed and attenuated streak development, with velocity fluctuations remaining poorly organised through a substantial portion of the domain. Whilst \gls{TEQ} and \gls{CPG} exhibit enhanced mode interactions leading to direct breakdown to turbulence, the \gls{TNE} case displays characteristic $\Lambda$-vortex structures consistent with observations in adiabatic hypersonic flows~\citep{Passiatore2021}, demonstrating the stabilising effect of vibrational nonequilibrium on secondary instability growth and coherent structure formation.

\begin{figure}[t]
    \centering
    \includegraphics[width=18cm,angle=-90]{resources/figures/breakdown/result_contour_yplus1_uvelocity.pdf}
    \caption[Instantaneous streamwise velocity ($u/u_e$) in $x-z$ plane during boundary layer transition for CPG, TEQ, and TNE models.]%
    {Instantaneous streamwise velocity ($u/u_e$) during boundary layer transition for (a) \gls{CPG}, (b) \gls{TEQ}, and (c) \gls{TNE} models.}
    \label{figure:velocity_snapshot}
\end{figure}

Instantaneous streamwise velocity snapshots in the $x$--$z$ plane at $y^+ = 1$ in the turbulent region are shown in \Cref{figure:velocity_snapshot}, revealing streak-like structures in all three cases with Mack mode interaction visible throughout the transitional region. For \gls{CPG} and \gls{TEQ}, modes develop earlier and reach significant amplitude at comparable streamwise locations, driving the observed wall-quantity overshoot~\citep{SchlatturEtAl2008}. The \gls{TNE} case exhibits delayed mode development at more advanced streamwise locations. Despite these spatial delays, vibrational nonequilibrium modulates the spatial progression of transition without fundamentally altering the underlying mechanisms of mode interaction and secondary instability growth. Coherent structures observed here are consistent with incompressible transition dynamics~\citep{SchoppaHussain2002,SchlatturEtAl2008} and with recent analysis of high-amplitude streak dynamics in hypersonic boundary layers~\citep{CaillaurdEtAl2025}, which provide critical insights into the modal interactions governing initial transition stages and nonlinear pathways to turbulence.


\subsubsection{Boundary layer evolution}

The development of the boundary layer during transition provides insight into changes in flow structure and momentum redistribution as instability modes grow and break down. Integral quantities such as the displacement thickness $\delta^*$ and momentum thickness $\theta^*$ characterise the overall profile growth and momentum deficit. \Cref{figure:dns_thermal_blgrowth} shows the evolution of these quantities for the three thermal models, highlighting how different thermochemical assumptions influence boundary-layer thickening rates and the spatial progression of transition.

As second-mode resonance develops, growth rates diverge markedly from the laminar regime. In the laminar and early linear instability stages, where vibrational excitation remains modest, \gls{CPG} and \gls{TNE} predictions exhibit close agreement, reflecting the secondary importance of thermal nonequilibrium effects. However, as transition progresses into the nonlinear breakdown region, a pronounced shift emerges: the \gls{CPG} model transitions towards behaviour more closely aligned with \gls{TEQ}, demonstrating accelerated displacement thickness growth. In contrast, the \gls{TNE} case manifests substantially attenuated thickness development, with thermal relaxation effects increasingly suppressing second-mode growth as vibrational nonequilibrium becomes dynamically significant. By the end of transition ($x/\delta^*_0 > 500$), all three cases undergo rapid boundary-layer thickening attributable to coherent structure generation and turbulent cascade processes. The shape factor $H = \delta^* / \theta$ in \Cref{table:global_quantities_models} exhibits a systematic reduction  from approximately $13$--$14$ in the laminar state to below $11$ in fully developed turbulence, reflecting progressive velocity-profile steepening as organised structures develop and dissipate.

\begin{figure}[t]
    \centering
    \includegraphics[width=11cm]{resources/figures/breakdown/result_dns_thermal_boundarylayergrowth.pdf}
    \caption[{Time- and spanwise-averaged boundary-layer growth along the streamwise direction for the DNS cases showing displacement thickness ($\delta^*$), momentum thickness ($\theta^*$), and enthalpy thickness ($\delta_h$) for CPG, TEQ, and TNE simulations.}]{Time- and spanwise-averaged boundary-layer growth along the streamwise direction for the DNS cases showing displacement thickness ($\delta^*$), momentum thickness ($\theta^*$), and enthalpy thickness ($\delta_h$) for \gls{CPG}, \gls{TEQ}, and \gls{TNE} simulations.}
    \label{figure:dns_thermal_blgrowth}
\end{figure}

This thickening of displacement and momentum thicknesses accompanies simultaneous reorganisation of the thermal field. The thermal boundary-layer evolution, illustrated through instantaneous temperature contours in \Cref{figure:thermal_snapshots}, reveals how energy redistribution accompanies momentum-layer development across the three models. From laminar through breakdown stages, temperature contours remain smooth and monotonic in the laminar region across all three models, with thermal energy confined to a thin layer adjacent to the wall. As second-mode resonance develops, incipient spanwise temperature modulations emerge, with \gls{TEQ} exhibiting more pronounced thermal striations than \gls{CPG} and \gls{TNE}. During the secondary-instability phase, spanwise temperature variations intensify dramatically, with coherent high-temperature streaks penetrating further into the outer boundary layer. By the nonlinear breakdown stage, complex three-dimensional thermal structures dominate, characterised by intermittent high-temperature filaments and low-temperature pockets reflecting the underlying quasi-streamwise vortex organisation. \gls{TEQ} displays the most vigorous thermal streaking consistent with accelerated transition, whilst \gls{TNE} shows delayed and more attenuated thermal organisation with weaker spanwise modulations, demonstrating that molecular energy redistribution fundamentally alters both momentum and thermal-structure development during hypersonic boundary-layer transition.

% \newpage
\begin{figure}[t]
    \centering
    \includegraphics[width=13.0cm,angle=0]{resources/figures/breakdown/result_breadkdown_thermal_Tslices.pdf}
    \caption{Instantaneous temperature ($T/T_e$) contours in the $x$--$y$ plane at a various spanwise locations during nonlinear breakdown and secondary-instability phases for (a) CPG, (b) TEQ, and (c) TNE models. Rows denote progression through the transition region.}
    \label{figure:thermal_snapshots}
\end{figure}
% \newpage



\begin{figure}[t]
    \centering
    \includegraphics[width=11.3cm]{resources/figures/breakdown/result_dns_thermal_fft.pdf}
    \caption{Normalised power spectral density and pressure rms distribution across the three thermochemical models. Panel (a) shows the forced-frequency response and panel (b) shows the spatial evolution of pressure fluctuation intensity during transition.}
    \label{figure:thermal_fft}
\end{figure}

\subsubsection{Frequency content}

FFT analysis was performed on wall pressure fluctuations in the $x$--$z$ plane during statistically steady flow phases across all three models to understand the effects of Mack modes in the transition process. The frequency spectra shown in \Cref{figure:thermal_fft}(a) reveal the response of each thermal model to the dual-frequency excitation. In the laminar region, pressure fluctuations remain minimal with power concentrated at low frequencies. As transition progresses and second-mode resonance develops, distinct spectral peaks emerge at the forced frequencies, with relative amplitudes indicating preferential mode amplification by boundary-layer stability characteristics. Marked spectral broadening is observed compared to the linear predictions of \Cref{chapter:disturbance_growth}, indicating nonlinear energy transfer driving coherent-structure breakdown.

The pressure rms distribution shown in \Cref{figure:thermal_fft}(b) shows significant differences across the three models. During resonance growth, \gls{CPG} exhibits the highest pressure rms intensity, whilst \gls{TEQ} and \gls{TNE} display comparatively attenuated fluctuations. However, during secondary instability and breakdown, pressure rms intensity increases substantially in the \gls{TNE} case, approaching or exceeding that of \gls{CPG}. This inversion of the rms hierarchy demonstrates the eventually delayed but vigorous nonlinear amplification characteristic of thermally nonequilibrium flows.

\begin{figure}[t]
    \centering
    \includegraphics[width=9.3cm]{resources/figures/breakdown/result_dns_thermal_pod_energy.pdf}
    \caption{Eigenvalue spectrum and cumulative energy content of the first 40 POD modes extracted from wall pressure fluctuations, demonstrating the concentration of modal energy in dominant coherent structures across CPG, TEQ, and TNE models.}
    \label{figure:pod_energy_variation}
\end{figure}

\subsubsection{Modal decomposition}

\begin{figure}[t]
    \centering
    \includegraphics[width=18.0cm,angle=-90]{resources/figures/breakdown/result_dns_thermal_contour_pod_mode1.pdf}
    \caption{Spatial structure of dominant POD mode extracted from wall pressure fluctuations during boundary-layer transition for (a) CPG, (b) TEQ, and (c) TNE models.}
    \label{figure:contour_podmode1}
\end{figure}


\gls{POD} provides a systematic approach to extract dominant coherent structures from flow fields by decomposing a field into orthogonal spatial modes ranked by energy content. The method decomposes a field $\psi(\mathbf{x},t)$ into a sum of spatial modes $\phi_n(\mathbf{x})$ with corresponding temporal coefficients $a_n(t)$, enabling identification of the most energetically significant structures governing flow dynamics. For incompressible flows, snapshot \gls{POD} has become a standard tool in boundary-layer transition studies, revealing the spatial organisation of large-scale coherent motions~\citep{Sirovich1987}. Spectral Proper Orthogonal Decomposition (SPOD) extends this framework by decomposing the flow into frequency-dependent modes, isolating coherent structures at specific frequencies and separating them from broadband acoustic noise \citep{TowneEtAl2018}. However, SPOD requires substantial computational resources and extensive data archives, making snapshot \gls{POD} more practical for \gls{DNS} studies with large datasets. Here, snapshot \gls{POD} is applied to wall pressure fluctuations in the $x$--$z$ plane to characterise the spatial structure and energy distribution of coherent pressure patterns during transition across the three thermochemical models. \gls{POD} was performed on wall pressure slices extracted in the $x$--$z$ plane at the wall, recorded over one flow-through time at intervals of 100 iterations. A standard singular value decomposition was employed to compute the spatial modes and their corresponding temporal coefficients, isolating the dominant pressure structures driving transition whilst filtering transient acoustic noise inherent to compressible simulations.

The spatial structure of the leading \gls{POD} mode in \Cref{figure:contour_podmode1} shows that coherent structures exhibit similar organisational characteristics across all three thermal models, with spanwise striations present in \gls{CPG}, \gls{TEQ}, and \gls{TNE}. In \gls{TNE}, however, the transition zone is more spatially compressed in the streamwise direction, concentrating modal content over a narrower region. Despite this, the transition is gradual, with prolonged energy transfer to higher modes and a slow rise in skin friction, as seen in \Cref{figure:mach_cf_st}. In contrast, \gls{TEQ} exhibits both an extended primary modal structure and the highest energy in higher modes, resulting in a sharper transition in skin friction. These observations indicate that thermal nonequilibrium fundamentally alters the relationship between primary instability spatial extent, modal energy cascade, and wall transition characteristics, reflecting qualitatively distinct transition pathways across the three models.

% The spatial structure of the leading \gls{POD} mode seen in \Cref{figure:contour_podmode1} reveals that coherent structures exhibit similar organisational characteristics across all three thermal models, with spanwise striations present in \gls{CPG}, \gls{TEQ}, and \gls{TNE}. However, the transition zone in \gls{TNE} is notably more spatially compressed in the streamwise direction, with transition-related modal content concentrated over a narrower region compared to the more extended transition zones in \gls{CPG} and \gls{TEQ}.  When considered alongside skin friction evolution in \Cref{figure:mach_cf_st}, which shows \gls{TEQ} and \gls{CPG} exhibiting sharp transitions with steep gradients whilst \gls{TNE} transitions more gradually over an extended streamwise range despite delayed onset, a complex picture emerges.

% The spatial compression of the forced Mack mode in \gls{TNE} does not correspond to abrupt transition. Rather, \gls{TNE}'s concentrated modal content is followed by prolonged energy transfer to higher modes and gradual skin friction rise. In contrast, \gls{TEQ} displays both an extended primary modal structure and the highest energy in higher modes, yet exhibits the sharpest transition in skin friction. These observations indicate that the relationship between primary instability spatial extent, modal energy cascade, and wall transition characteristics is fundamentally altered by thermal nonequilibrium, reflecting qualitatively distinct transition pathways across the three models. The sharp transition in \gls{TEQ} may reflect rapid nonlinear saturation driven by accumulated linear instability growth coupled with wall-imposed forcing, whilst the gradual transition in \gls{TNE} indicates that vibrational nonequilibrium modifies nonlinear energy-transfer mechanisms.



\subsubsection{Disturbance growth}

The evolution of disturbance kinetic energy across the transition process was examined to understand flow behaviour. Disturbance kinetic energy is defined as
\begin{equation}\label{equation:tke}
e_{dke} = \frac{1}{2}\left(\widetilde{u}'^2 + \widetilde{v}'^2 + \widetilde{w}'^2\right),
\end{equation}
where $\widetilde{u}'$, $\widetilde{v}'$, and $\widetilde{w}'$ are the root-mean-square velocity fluctuations obtained from Favre-averaged variables in the streamwise, wall-normal, and spanwise directions, respectively. This quantity directly quantifies the energy contained in velocity perturbations and reveals how instability modes amplify and redistribute energy across spatial regions during transition. \Cref{figure:eke_budget} presents contours of disturbance kinetic energy in the $x$--$y$ plane with time and spanwise averaging, for the Mach 6, $T_w = 1800~\text{K}$ case across \gls{CPG}, \gls{TEQ}, and \gls{TNE} formulations. The aspect ratio is intentionally non-unity to clearly visualise the streamwise evolution of disturbance energy as it concentrates near the wall during linear growth, intensifies during secondary instability, and eventually extends throughout the boundary layer as turbulence develops. 

\begin{figure}[t]
    \centering
    \fontsize{9}{10}\selectfont
    \includegraphics[width=0.80\textwidth]{resources/figures/breakdown/result_dns_thermal_tke_m6tw18.pdf}
    \caption{Time- and spanwise-averaged disturbance kinetic energy contours in the $x$--$y$ plane for Mach 6, $T_w = 1800~\text{K}$ comparing CPG, TEQ, and TNE formulations.}
    \label{figure:eke_budget}
\end{figure}

The disturbance kinetic energy contours presented in \Cref{figure:eke_budget} provide direct quantification of energy evolution across the transition process. For \gls{CPG} and \gls{TEQ}, the spatial extent of energy is distributed across the wall-normal direction early in transition, with distinct peaks occurring at comparable streamwise locations and reflecting rapid transition onset. In contrast, \gls{TNE} exhibits more confined, intense energy concentrations near the wall initially, which extend progressively into the boundary layer at downstream locations. This spatial confinement reflects the delayed amplification of disturbances, yet the ultimately higher energy magnitude in the \gls{TNE} case demonstrates the vigorous nonlinear growth following secondary instability emergence. The wall-normal penetration of disturbance energy differs markedly in the fully developed region, with \gls{TNE} energy extending higher relative to boundary-layer thickness, indicating fundamentally different turbulent structure organisation modulated by vibrational nonequilibrium. The earlier saturation of energy growth in \gls{CPG} and \gls{TEQ} contrasts with the sustained energy rise in \gls{TNE}, reflecting the prolonged transition pathway in thermally nonequilibrium flows. These observations unify the transition behaviour revealed by skin friction, spectral content, modal decomposition, and mean-flow statistics, demonstrating that thermal nonequilibrium systematically delays transition whilst preserving the fundamental instability mechanisms.

\subsection{Turbulent flow}

Once the boundary layer transitions through nonlinear breakdown and secondary-instability phases, thermal structures progressively dissipate into small-scale turbulent fluctuations whilst mean-flow properties asymptotically approach turbulent equilibrium values. In this fully turbulent regime, the three thermochemical models exhibit markedly different turbulence characteristics despite comparable integral boundary-layer growth. To isolate and quantify model-dependent effects on turbulence, velocity statistics and Reynolds stress distributions are examined. Statistical averaging was performed over $1.5$ flow-through times following establishment of fully turbulent conditions, corresponding to approximately 50 eddy turnover times~\citep{Pirozzoli2004}, sufficient to capture converged statistics whilst resolving the high-frequency acoustic modes characteristic of hypersonic boundary layers. Representative stations within the turbulent region are analysed for each model, with streamwise locations specified in \Cref{table:global_quantities_models}.

\subsubsection{Velocity and temperature statistics}
A detailed assessment of velocity statistics is essential for characterising the turbulent boundary layer. Mean velocity profiles provide direct information on the redistribution of momentum by turbulence and form the basis for testing scaling laws and transformation approaches. Higher-order statistics, such as Reynolds stresses, quantify the intensity and anisotropy of turbulent fluctuations and reveal how energy is transferred amongst flow components. Together, these measures offer a comprehensive view of the turbulent structure and allow differences between thermal models to be identified and interpreted.

In compressible boundary layers, direct comparison with incompressible velocity profiles requires a transformation that accounts for density and viscosity variations near the wall. The van Driest transformation \citep{VanDriest1952} extends the original scaling to include density effects, offering improved collapse of velocity profiles in the inner region. This provides a common reference for assessing compressible turbulent boundary layers and comparing different thermal models on a consistent basis. The mean velocity and wall-normal coordinate $y$ are scaled using the friction velocity $u_{\tau}$ and the viscous length scale $\nu_w/u_{\tau}$. The transformed velocity is defined as
\begin{equation}
    u^{+} = \int_{0}^{\bar{u}} \sqrt{\frac{\bar{\rho}}{\rho_w}} \, \mathrm{d}u ,
\end{equation}
where $\bar{\rho}$ is the time-averaged local mean density and $\rho_w$ is the density at the wall.

The resulting mean velocity profile should collapse onto the universal law of the wall, matching the viscous sublayer ($u^{+}=y^{+}$ for $y^{+}<10$) and the logarithmic region at larger wall-normal distances:
\begin{equation}
u^{+} = \frac{1}{\kappa}\ln(y^{+}) + C, 
\qquad \kappa = 0.41, \quad C = 5.2.
\end{equation}

Recent developments have introduced improved velocity transformations for compressible profiles. Semilocal scaling concepts~\citep{TrettelLarsson2016,GriffinEtAl2021,KumarLarsson2022} emphasise the role of variable density and viscosity in defining an effective viscous length scale. Investigations of van Driest transformations accounting for intrinsic compressibility~\citep{HasanEtal2023,HasanEtal2024} have revealed that this effect is a characteristic feature of thermochemical nonequilibrium turbulence~\citep{Williamsetal2025}, whilst cold-wall conditions suppress it~\citep{Zhangetal2018}. The present work employs the transformations suggested by~\citep{HasanEtal2024} denoted as $u^{+}_{H}$, which provide enhanced accuracy across diverse wall conditions and Mach numbers.

Under the assumption of universality of the turbulent shear stress in the near-wall region, the semilocal wall coordinate is related to the wall-normal coordinate through
\begin{equation}
y^{*} = \frac{\delta_{v}}{\delta^{*}_{v}}\, y^{+},
\end{equation} where $\delta_v$ where, $\delta_v$ is the viscous length scale based on local wall properties, $\delta_v = \nu_w/u_\tau$, while $\delta_v^*$ is a reference viscous length scale, typically based on freestream or averaged properties, used to account for variable density and viscosity effects in the semilocal scaling.


\begin{figure}[t]
    \centering
    \fontsize{9}{10}\selectfont
    \includegraphics[width=0.60\textwidth]{resources/figures/breakdown/result_dns_velocity_profiles.pdf}
    \caption{Profiles of velocity and temperature statistics plotted against the semilocal wall coordinate $y^{*}$, obtained from the van Driest transformation, for Mach 6, $T_w=1800$ K comparing CPG, TEQ, and TNE formulations.}
    \label{figure:velocity_statistics}
\end{figure}

The velocity profiles collapse well onto the expected scaling laws when expressed in terms of the semilocal wall coordinate $y^{*}$, following the linear relation in the viscous sublayer and approaching the logarithmic law in the overlap region as shown in \Cref{figure:velocity_statistics}. This confirms that the van Driest transformation provides a consistent framework across the different thermal models. The turbulent Mach number $M_t$ exhibits clear thermal sensitivity, with \gls{TEQ} producing the highest peak, followed by \gls{TNE}, and \gls{CPG} remaining lowest. This ranking mirrors observations in adiabatic flows~\citep{Duan2011}, where increasing freestream Mach number produces similar shifts in $M_t$ profiles. These results demonstrate that vibrational nonequilibrium fundamentally alters turbulent structure, with vibrational excitation effects rivalling those of bulk flow acceleration. Despite elevated disturbance kinetic energy in \gls{TNE}, the turbulent Mach number remains intermediate, indicating that vibrational relaxation redistributes energy towards non-acoustic fluctuation modes rather than enhancing acoustic velocity fluctuations comparable to \gls{TEQ}.



\subsubsection{Favre-averaged turbulence stresses}

The turbulence stress components at $x/\delta^*_0 = 654$ and $740$, extracted from Favre-averaged statistics and normalised by the mean wall shear stress $\tau_w$, exhibit the canonical ordering $\overline{\rho}\,\widetilde{u''u''} > \overline{\rho}\,\widetilde{v''v''} \sim \overline{\rho}\,\widetilde{w''w''}$ as shown in \Cref{figure:thermal_reynolds_stresses}. The streamwise component reaches $\overline{\rho}\,\widetilde{u''u''}/\tau_w \approx 10.5$ at $y^* \approx 30$, whilst $\overline{\rho}\,\widetilde{v''v''}$ and $\overline{\rho}\,\widetilde{w''w''}$ attain peak values of approximately $3.5$ and $1.0$--$1.3$ respectively, consistent with canonical DNS of hypersonic turbulent boundary layers~\citep{ZhangDuanChoudhari2018,XuEtal2021}.

Critically, across all stress components at both stations, the \gls{CPG}, \gls{TEQ}, and \gls{TNE} closures collapse onto nearly indistinguishable profiles. This collapse demonstrates that once turbulence is fully developed, vibrational nonequilibrium produces negligible changes to Reynolds stress structure. The result reflects differing transition Reynolds numbers across the three thermal models rather than any fundamental alteration to turbulent dynamics. These findings are consistent with recent studies~\citep{DiRenzo2021,Williamsetal2025}, confirming that turbulence drives nonequilibrium effects rather than the reverse, with thermal excitations important only for the transition process.


\begin{figure}[t]
    \centering
    \fontsize{9}{10}\selectfont
    \includegraphics[width=0.80\textwidth]{resources/figures/breakdown/result_dns_thermal_reynolds_stresses.pdf}
    \caption{Favre-Reynolds stress components and root-mean-square fluctuations plotted against semilocal wall coordinate $y^*$ for Mach 6, $T_w=1800$ K. Left column shows $\overline{\rho}\,\widetilde{u''u''}$, $\overline{\rho}\,\widetilde{v''v''}$, $\overline{\rho}\,\widetilde{w''w''}$, and $-\overline{\rho}\,\widetilde{u''v''}$ normalised by wall shear stress. Right column shows corresponding rms velocity fluctuations. Comparisons between CPG, TEQ, and TNE formulations at multiple streamwise stations.}
    \label{figure:thermal_reynolds_stresses}
\end{figure}


\begin{figure}[t]
    \centering
    \fontsize{9}{10}\selectfont
    \includegraphics[width=0.80\textwidth]{resources/figures/breakdown/result_dns_thermal_rhop_rhoT.pdf}
    \caption{Profiles of Favre-averaged pressure (top row) and temperature (bottom row) with corresponding normalised root-mean-square fluctuations plotted against semilocal wall coordinate $y^*$ for Mach 6, $T_w=1800$ K. Left column shows mean profiles, with right showing the rms fluctuations. Legend: CPG (\protect\cpgline), TEQ (\protect\teqline), TNE (\protect\tneline).}
    \label{figure:thermal_rho_T}
\end{figure}

The mean pressure and temperature profiles, along with their rms fluctuations, are presented across the boundary layer in \Cref{figure:thermal_rho_T}. The mean pressure varies only weakly in the wall-normal direction, with all three formulations remaining close to their edge values. Small differences appear in the recovery region, where \gls{TNE} relaxes more gradually than \gls{CPG} and \gls{TEQ}, reflecting slower adjustment associated with vibrational energy modes. The mean temperature shows similar behaviour, with \gls{TNE} remaining elevated further into the boundary layer. The fluctuation levels display more distinct separation, with \gls{TNE} producing larger rms pressure and temperature fluctuations than both \gls{CPG} and \gls{TEQ}. The additional degrees of freedom introduced by vibrational nonequilibrium strengthen coupling between thermal and momentum fields, increasing compressible-mode activity and enhancing turbulent mixing. Consequently, both pressure and temperature variances remain higher throughout the log layer and outer region.

The temperature results exhibit particularly slow relaxation. Whilst linear DNS studies suggest downstream \gls{TNE} temperature should approach \gls{TEQ} more closely than \gls{CPG}, fully turbulent simulation shows markedly different behaviour. Turbulence acts to delay thermal equilibration, with \gls{TNE} mean temperature remaining clearly separated from \gls{TEQ} well into the boundary layer. These observations indicate that velocity statistics adjust rapidly to classical turbulent behaviour whilst thermal nonequilibrium persists much further downstream. Vibrational nonequilibrium is promoted by turbulence~\citep{FievetEtAl2019}, meaning that once the flow transitions, it tends to remain nonequilibrium over extended distances. This underscores that vibrational nonequilibrium effects are particularly significant during transition when the flow remains far from equilibrium.


\section{Comparison of DNS and ILES cases}
\label{section:gridstudy}

The fine-resolution \gls{DNS} grids presented in the preceding sections comprised approximately $5200 \times 600 \times 275$ nodes in the streamwise, wall-normal, and spanwise directions respectively, with near-wall spacing designed to resolve viscous sublayer dynamics and instability modes. The domain dimensions are $800\delta^*_0 \times 100\delta^*_0 \times 20\delta^*_0$ in the streamwise, wall-normal, and spanwise directions, yielding grid spacings of $\Delta x^+ \approx 10.66$, $\Delta y^+_w \approx 0.94$, and $\Delta z^+ \approx 4.95$ respectively.

\begin{figure}[t]
    \centering
    \fontsize{9}{10}\selectfont
    \includegraphics[width=0.88\textwidth]{resources/figures/breakdown/result_wallquantities_cf_st_grid.pdf}
    \caption{Streamwise evolution of skin friction coefficient $C_f$ and Stanton number $St$ for the fine-resolution DNS cases. Panels show (a) CPG, (b) TEQ, and (c) TNE formulations with laminar and turbulent reference solutions. Solid lines represent DNS, markers represent ILES.}
    \label{figure:gridstudy_cfst}
\end{figure}

In addition to the fine-resolution \gls{DNS} cases, a coarser grid strategy was adopted to evaluate \gls{ILES} as a lower-cost alternative for transition studies. \gls{ILES} relies on the dissipative properties of high-order numerical schemes to implicitly model subgrid-scale turbulence without explicit subgrid models~\citep{PoggieEtAl2015}. The \gls{ILES} grids comprised approximately $2300 \times 500 \times 165$ nodes over the same domain dimensions. This yields streamwise and spanwise resolutions of $\Delta x^+ \approx 23.33$ and $\Delta z^+ \approx 12.12$ respectively. This coarser resolution substantially reduces computational cost, with \gls{ILES} simulations approximately 6.5 times less expensive than \gls{DNS} due to reduced grid requirements, whilst capturing the dominant instability mechanisms.

\begin{figure}[t]
    \centering
    \fontsize{9}{10}\selectfont
    \includegraphics[width=0.88\textwidth]{resources/figures/breakdown/result_dns_thermal_boundarylayergrowth_grid3x3.pdf}
    \caption{Streamwise evolution of boundary-layer growth showing displacement thickness $\delta^*$, momentum thickness $\theta^*$, and enthalpy thickness $\delta_h$ for (a) CPG, (b) TEQ, and (c) TNE formulations. Solid lines represent DNS, markers represent ILES.}
    \label{figure:gridstudy_boundarylayergrowth}
\end{figure}

\begin{table}[htb]
\centering
\caption{Grid resolution parameters for DNS and ILES cases at Mach~6, wall temperature $T_w = 1800$~K.}
\label{table:grid_parameters}
\begin{tabular}{llcccccc}
\toprule
\textbf{Model} & \textbf{Case} & $\Delta x^+$ & $\Delta y^+_w$ & $\Delta y^+_{\delta_{99}}$ & $\Delta z^+$ & $L_z^+$ \\
\midrule
\multirow{2}{*}{CPG} & DNS & 10.66 & 0.94 & 9.15 & 4.95 & 1385 \\
& ILES & 24.18 & 0.88 & 8.52 & 16.24 & 1546 \\
\midrule
\multirow{2}{*}{TEQ} & DNS & 8.35 & 0.73 & 6.89 & 3.88 & 1085 \\
& ILES & 23.46 & 0.92 & 7.18 & 14.67 & 1239 \\
\midrule
\multirow{2}{*}{TNE} & DNS & 10.97 & 0.96 & 8.27 & 5.09 & 1426 \\
& ILES & 26.52 & 0.84 & 9.74 & 17.43 & 1678 \\
\bottomrule
\end{tabular}
\end{table}



To assess the suitability of \gls{ILES} for the present transition study, comparisons were made between corresponding \gls{DNS} and \gls{ILES} cases across the three thermal models. \Cref{figure:gridstudy_cfst} presents the streamwise evolution of skin friction coefficient $C_f$ and Stanton number $St$. The location of the Mack-mode resonance remains consistent between \gls{DNS} and \gls{ILES}, and the transition onset location is preserved. Notably, \gls{TNE} exhibits the sharpest transition in skin friction, reaching the transition threshold significantly earlier than \gls{CPG} and \gls{TEQ}. The primary differences between \gls{DNS} and \gls{ILES} manifest in the turbulent intensity and the sharpness of transition, where \gls{ILES} exhibits slightly smoother gradients in the transition region. This behaviour is consistent with observations that WENO and TENO schemes produce comparable results between \gls{DNS} and \gls{ILES}~\citep{HamzehloEtAl2021}.



\Cref{figure:gridstudy_boundarylayergrowth} shows the streamwise evolution of boundary-layer integral parameters. The displacement and momentum thicknesses track closely between \gls{DNS} and \gls{ILES} across all three thermal models, with \gls{TEQ} showing the tightest agreement and \gls{TNE} demonstrating remarkably consistent behaviour and trends throughout the domain. This close correspondence extends to the enthalpy thickness, indicating that \gls{ILES} accurately captures the overall boundary-layer growth and the transition location across thermal models.

\begin{figure}[t]
    \centering
    \fontsize{9}{10}\selectfont
    \includegraphics[width=0.88\textwidth]{resources/figures/breakdown/result_dns_grid_study_upH_Mt_k.pdf}
    \caption{Streamwise evolution of wall-normal velocity profiles $u^+$, turbulent Mach number $M_t$, and turbulent kinetic energy $k$ for (a) CPG, (b) TEQ, and (c) TNE formulations. Solid lines represent DNS, markers represent ILES.}
    \label{figure:grid_study_upH_Mt_k}
\end{figure}

In the fully turbulent region, \Cref{figure:grid_study_upH_Mt_k} reveals minor differences in turbulence statistics. The \gls{ILES} velocity profiles $u^+$ follow the viscous sublayer and log law closely with \gls{DNS}, but begin to diverge in the buffer region and continue diverging throughout the log layer. The turbulent Mach number $M_t$ and turbulent kinetic energy $k$ (defined analogously to the disturbance kinetic energy in equation \Cref{equation:tke}
) show similarly modest overprediction. These differences arise from the coarser grid resolution underresolving fine-scale coherent structures in the turbulent boundary layer. However, the overall agreement is remarkably good considering the substantial reduction in computational cost, with \gls{ILES} approximately 6.4 times less expensive than \gls{DNS}. Consequently, whilst \gls{ILES} shows small errors in fully turbulent statistics, it provides a cost-effective pathway for systematic comparisons of transition location and early transition dynamics across multiple thermal models and flow conditions.

\section{Role of Mach number and wall temperature}

The combined influence of Mach number and wall temperature on stability and transition in hypersonic boundary layers is assessed within vibrational nonequilibrium. Increasing Mach number amplifies compressibility effects that dampen acoustic-trapped disturbances~\citep{BitterShepherd2015}, whilst wall temperature variation modifies the mean thermal gradients driving instability growth. Simulations are performed for \gls{TNE} flows at Mach 6 and Mach 8 with cold wall $T_w = 1200~\text{K}$ and hot wall $T_w = 1800~\text{K}$ conditions using \gls{ILES}. The computational domains span $800\delta^*_0$ for Mach 6 and $1600\delta^*_0$ for Mach 8 to accommodate different instability wavelengths and growth rates. The streamwise spacing was adjusted to maintain consistent wall-normal and spanwise resolution across Mach numbers.

Disturbance frequencies are imposed as $\hat{f} = 0.15$ and $0.16$ for Mach 6, and $\hat{f} = 0.10$ and $0.11$ for Mach 8, selected to capture the resonance peak region at $x/\delta^*_0 \approx 200$--$400$ where modal growth is most pronounced and trigger transition within the domain. The $N$-factor analysis in \Cref{chapter:disturbance_growth} reveals that the Mach 6 cold wall case exhibits the highest amplification, with the hot wall showing suppression, whilst both Mach 8 cases display further stabilisation. This demonstrates that increased Mach number and elevated wall temperature both stabilise modal growth. Vibrational nonequilibrium effects persist across all Mach numbers and wall temperatures examined.


\begin{figure}[t]
    \centering
    \fontsize{9}{10}\selectfont
    \includegraphics[width=0.95\textwidth]{resources/figures/breakdown/result_wallquantities_cf_st_tne.pdf}
    \caption{Streamwise evolution of skin friction coefficient $C_f$ and Stanton number $St$ for TNE formulations showing (a) Mach 6, $T_w=1200$ K; (b) Mach 6, $T_w=1800$ K; (c) Mach 8, $T_w=1200$ K; and (d) Mach 8, $T_w=1800$ K with laminar and turbulent reference solutions.}
    \label{figure:mach_cf_st}
\end{figure}


\begin{figure}[h]
    \centering
    \fontsize{9}{10}\selectfont
    \includegraphics[width=0.75\textwidth]{resources/figures/breakdown/result_wallquantities_cq_tne.pdf}
    \caption{Streamwise evolution of translational--rotational ($C_{q,\mathrm{tr}}$, black) and vibrational ($C_{q,\mathrm{v}}$, red, scaled $\times 10$) heat flux coefficients for TNE formulations showing (a) Mach 6, $T_w=1200$ K; (b) Mach 6, $T_w=1800$ K; (c) Mach 8, $T_w=1200$ K; and (d) Mach 8, $T_w=1800$ K.}
    \label{figure:cqplot}
\end{figure}

\subsection{Global mean flow and surface properties}

The skin friction coefficient $C_f$, Stanton number $St$, and normalised heat transfer coefficient $C_{q,w} = q_w / (\rho_e u_e^3)$ provide quantitative measures of wall shear stress and surface heating across all Mach numbers and thermal conditions~\citep{Volpiani2021}. \Cref{figure:mach_cf_st} presents the skin friction and Stanton number evolution across all cases. A pronounced overshoot in $C_f$ occurs upon transition onset across all cases. The resonance peak is pronounced in the Mach 6 cases, consistent with earlier DNS studies, but considerably subdued in the Mach 8 cases. This attenuation reflects the lower $N$-factor response at higher Mach number, with maximum $N$-factors of 8 and 6 for Mach 6 cases compared to 4 and 2 for Mach 8 cases. The transition sequence follows the expected progression, with the Mach 6 cold wall case transitioning earliest and the Mach 8 hot wall case exhibiting the most extended transition process. Notably, despite employing the same grid resolution as the Mach 6 cases, the Mach 8 simulations exhibit a DNS-level spike in skin friction during transition, indicating that higher compressibility effects at elevated Mach number produce sharper transition characteristics even on coarser grids. The Stanton number follows similar trends to skin friction, with pronounced peaks in the Mach 6 cases and more gradual evolution in the Mach 8 cases.

\begin{figure}[t]
    \centering
    \includegraphics[width=8.9cm]{resources/figures/breakdown/result_tne_amplitude_nfactor.pdf}
    \caption{Streamwise evolution of rms amplitude (left column) and cumulative $N$-factor (right column) for the prescribed forcing frequencies showing (a) Mach 6, $T_w=1200$ K ; (b) Mach 6, $T_w=1800$ K; (c) Mach 8, $T_w=1200$ K; and (d) Mach 8, $T_w=1800$ K  for TNE conditions.}
    \label{figure:tne_amplitude_nfactor}
\end{figure}

To assess the influence of vibrational nonequilibrium, the translational--rotational and vibrational heat flux components are examined. The net total heat flux is represented by the sum of $C_{q,\mathrm{tr}}$ and $C_{q,\mathrm{v}}$, normalised as $C_{q,w} = q_w / (\rho_e u_e^3)$ as shown in \Cref{figure:cqplot}, with $C_{q,\mathrm{v}}$ scaled by a factor of ten. The translational--rotational and vibrational heat flux components exhibit strikingly different behaviours across the four cases. In the Mach 6 cold wall case, $C_{q,\mathrm{v}}$ remains suppressed during the laminar and linear growth regions, rising sharply coincident with secondary instability and transition. The Mach 6 hot wall case shows elevated vibrational heat flux throughout, with significant levels in the laminar region that amplify further during transition. At Mach 8, the cold wall case exhibits minimal vibrational contribution across the entire domain, whilst the hot wall case shows weak but persistent vibrational activity that gradually increases during the transition process. This dependence on wall temperature reveals that vibrational excitation is promoted by elevated temperatures, with cold walls suppressing vibrational modes until nonlinear breakdown occurs.





\subsection{Modal growth}

One of the distinct features of the transition studies presented in this chapter has been the resonance peaks in the skin friction and Stanton number plot. The resonance peak was attributed to the initial growth of the 2D Mack mode~\citep{Franko2013}, and the peak was observed in subsequent studies of hypersonic boundary layers~\citep{DiRenzo2021}, though absent in others. Both observations stem from adiabatic wall conditions. The prescribed forcing frequencies are selected to excite the mode at its point of maximum linear amplification, as identified in the $N$-factor analysis presented in \Cref{figure:tne_amplitude_nfactor}, with the peak location corresponding to where the $N$-factor reaches its maximum. The resonance peak prominence is strongly modulated by thermal conditions, appearing pronounced in the Mach 6 cases where $N$-factors reach 8 and 6 respectively, but considerably subdued in the Mach 8 cases where $N$-factors remain at 4 and 2. Notably, despite employing identical forcing frequency amplitudes across all cases, the resonance peaks are absent in the Mach 8 simulations, reflecting the substantially lower linear growth rates at elevated Mach number that fail to produce sufficient wall-quantity amplification to manifest as distinct peaks. Cold wall conditions enhance the peak through increased thermal gradients driving instability growth, whilst elevated wall temperature and higher Mach number both suppress linear amplification. The resonance peaks transition smoothly into nonlinear breakdown as secondary instabilities emerge, reflecting the natural progression from linear to nonlinear dynamics.


\begin{figure}[t]
    \centering
    \includegraphics[width=18.3cm,angle=-90]{resources/figures/breakdown/result_iles_thermal_contour_pod_mode1.pdf}
    \caption{Streamwise evolution of the first POD of pressure fluctuations in the spanwise direction for TNE formulations showing (a) Mach 6, $T_w=1200$ K; (b) Mach 6, $T_w=1800$ K; (c) Mach 8, $T_w=1200$ K; and (d) Mach 8, $T_w=1800$ K.}
    \label{figure:iles_pod_mode1}
\end{figure}

To further isolate the dominant coherent structures driving transition across the parametric range, proper orthogonal decomposition was applied to pressure fluctuations in the spanwise direction. \Cref{figure:iles_pod_mode1} presents the first POD mode of pressure fluctuations across all four thermal conditions. The Mach 6 cases display diffuse, irregular patterns that rapidly break down into small-scale fluctuations, reflecting the vigorous nonlinear interactions characteristic of nonlinear growth. In contrast, the Mach 8 cases exhibit highly organised, periodic spanwise structures that persist throughout the modal amplification regime, indicating a strong modulation by the forcing frequency. The spanwise wavelengths evident in the \gls{POD} mode align with the imposed forcing wavelengths, confirming that the dominant pressure structures are directly driven by the prescribed disturbances.

The presence of the resonance peak in the Mach 6 cases and its absence in the Mach 8 cases reflects the competition between linear amplification and nonlinear saturation. At Mach 6, the higher linear growth rates allow the forced mode to reach substantial amplitude over a distinct spatial region before secondary instabilities dominate, producing the pronounced resonance peak. At Mach 8, the suppressed linear growth rates prevent the primary mode from achieving sufficient amplitude to generate a detectable peak, with transition driven instead by a more gradual energy cascade directly into nonlinear breakdown.

\subsection{Streak breakdown}

Streak breakdown via secondary instability is a critical transition mechanism in boundary layers. Streaks are quasi-streamwise vortical structures generated through the lift-up mechanism, which amplify until secondary instability triggers breakdown to turbulence. In incompressible flows, secondary instabilities were characterised using Floquet theory~\citep{Andersson2001}, revealing sinuous and varicose modes with critical amplitudes of $26\,\%$ and $37\,\%$ of freestream velocity respectively. ~\cite{BrandtBottaro2001} extended this work to include nonlinear saturation of streaks and full breakdown to turbulence through forced secondary instability. However, the critical amplitude threshold for compressible flows remains poorly established, with no dominant mechanism or commonly accepted streak amplitude percentage identified in high-speed simulations~\citep{Franko2013}.

Streak amplitude is computed from the velocity field as the maximum spanwise variation of streamwise velocity, normalised by the freestream velocity. For each streamwise location and wall-normal position, the spanwise deviation is calculated as
\begin{equation}
A_s(x) = \frac{\max_z u(x, y, z) - \min_z u(x, y, z)}{u_e},
\end{equation}
where the streak amplitude at each streamwise station is then defined as the maximum value across the wall-normal direction.

\begin{figure}[H]
    \centering
    \fontsize{9}{10}\selectfont
    \includegraphics[width=0.90\textwidth]{resources/figures/breakdown/result_iles_streak_amplitude.pdf}
    \caption{Streak amplitude growth in the streamwise direciton. Results are shown for TNE studies.}
    \label{figure:streakbreakdown}
\end{figure}


\begin{figure}[H]
    \centering
    \includegraphics[width=20.3cm,angle=-90]{resources/figures/breakdown/result_iles_streak_contour.pdf}
    \caption{Streamwise evolution of streak in $u/u_e$ slice in the spanwise direction at $y^{+}=1$ for TNE formulations showing (a) Mach 6, $T_w=1200$ K; (b) Mach 6, $T_w=1800$ K; (c) Mach 8, $T_w=1200$ K; and (d) Mach 8, $T_w=1800$ K. }
    \label{figure:iles_streak_amplitude}
\end{figure}

The development of sinusoidal streak structures as mentioned in \citep{Andersson2001} is evident across all thermal conditions shown in \Cref{figure:iles_streak_amplitude}. The Mach 6 cases exhibit rapid streak formation and amplification beginning in the secondary instability region, with coherent spanwise-periodic structures emerging around $x/\delta^*_0 \approx 300$ and growing to substantial amplitudes approaching the breakdown region. The Mach 8 cases show more gradual streak development with lower peak amplitudes, consistent with the weaker secondary instability growth at higher Mach number. Streak amplitude increases systematically through the transition zone, with the most vigorous growth occurring in the Mach 6 cold wall case where linear instability is most pronounced. The sinusoidal organisation reflects the forcing-induced modulation of the flow, demonstrating that secondary instabilities preferentially amplify structures aligned with the imposed spanwise wavelengths.




This metric quantifies the strength of quasi-streamwise vortical structures and provides a direct measure of secondary instability susceptibility. The streak amplitude evolution across the four thermal nonequilibrium cases is shown in \Cref{figure:streakbreakdown}. At Mach 6, $T_w = 1200$ K, rapid amplification occurs with the secondary instability threshold crossed early at approximately $32\,\%$ of freestream velocity. Increasing wall temperature to $1800$ K delays streak growth, reflecting the stabilising effect of elevated wall temperature on secondary instability. At Mach 8, streak amplitudes remain suppressed at both wall temperatures, indicating that compressibility and wall heating combine to inhibit secondary instability at higher Mach number. The thermal nonequilibrium effect manifests through altered thermodynamic properties governing secondary instability growth rates. Energy partitioning between translational and vibrational modes introduces a stabilising influence on streak amplification, consistent with the delayed transition observed in \gls{TNE} mean flow statistics. These results emphasise the importance of accounting for compressibility, wall effects, and thermal nonequilibrium when predicting transition in hypersonic flows.



\begin{figure}[t]
    \centering
    \fontsize{9}{10}\selectfont
    \includegraphics[width=0.80\textwidth]{resources/figures/breakdown/result_iles_contour_TmTv_fulldomain.pdf}
    \caption{Streamwise evolution of mean temperature field normalised by freestream conditions for TNE formulations showing (a) Mach 6, $T_w=1200$ K; (b) Mach 6, $T_w=1800$ K; (c) Mach 8, $T_w=1200$ K; and (d) Mach 8, $T_w=1800$ K.}
    \label{figure:contour_TmTv}
\end{figure}


\subsection{Thermal nonequilibrium effects}

In this section, vibrational nonequilibrium effects across the four thermal conditions are analysed to understand how Mach number and wall temperature modulate transition behaviour in hypersonic flows. Figure~\ref{figure:contour_TmTv} presents contours of time- and spanwise-averaged vibrational temperature deviation $\overline{T_v - T}$ normalised by freestream temperature throughout the domain, with non-unity aspect ratio to reveal regions of thermal nonequilibrium. The maximum vibrational nonequilibrium occurs within the transition region across all cases, where secondary instabilities drive rapid energy redistribution between translational and vibrational modes. In the fully turbulent region downstream, the nonequilibrium zones persist with reduced magnitude, indicating that vibrational relaxation effects become less dynamically significant once transition is complete. The spatial extent and intensity of nonequilibrium vary systematically with Mach number and wall temperature, with elevated Mach number and cold walls suppressing vibrational excitation whilst lower Mach number and hot walls promoting stronger thermal nonequilibrium effects.

\begin{figure}[t]
    \centering
    \fontsize{9}{10}\selectfont
    \includegraphics[width=0.80\textwidth]{resources/figures/breakdown/result_iles_Damkohler_number.pdf}
    \caption{Wall-normal profiles of Damk\"{o}hler number defined as the ratio of viscous timescale to vibrational relaxation timescale across five streamwise stations (laminar, resonance, secondary instability, transition, and turbulent regions) for TNE formulations showing (a) Mach 6, $T_w=1200$ K; (b) Mach 6, $T_w=1800$ K; (c) Mach 8, $T_w=1200$ K; and (d) Mach 8, $T_w=1800$ K. Colours denote streamwise position, with darker shades indicating downstream evolution.}
    \label{figure:damkohlernumber}
\end{figure}

\begin{figure}[t]
    \centering
    \fontsize{9}{10}\selectfont
    \includegraphics[width=0.80\textwidth]{resources/figures/breakdown/result_iles_thermal_TmTv.pdf}
    \caption{Wall-normal profiles of vibrational temperature deviation $(T_v - T)/T_\infty$ across five streamwise stations (laminar, resonance, secondary instability, transition, and turbulent regions) for TNE formulations showing (a) Mach 6, $T_w=1200$ K; (b) Mach 6, $T_w=1800$ K; (c) Mach 8, $T_w=1200$ K; and (d) Mach 8, $T_w=1800$ K. Left column presents linear wall-normal coordinates $y/\delta^{*}$, whilst right column presents semi-logarithmic coordinates to reveal inner-layer structure. Colours denote streamwise position, with darker shades indicating downstream evolution.}
    \label{figure:iles_tmtv_profiles}
\end{figure}


These observations are consistent with the findings of~\cite{FievetEtAl2019} and~\cite{Passiatore2021}, who identified that over-excited vibrational states develop in regions where turbulent intensity decreases whilst vibrational energy relaxes slowly. However, the present results show that the transition process itself drives the flow into thermal nonequilibrium through the intense energy redistribution accompanying secondary instability growth. The concentration of nonequilibrium in the transition region demonstrates that the instability-driven acceleration and deceleration of fluid parcels creates thermodynamic conditions where the vibrational relaxation timescale becomes comparable to the hydrodynamic timescale, preventing rapid equilibration and generating over-excited vibrational states. To quantify this balance and understand how the competing timescales determine the strength of TNE effects across the parametric range examined, a Damköhler number analysis based on the viscous timescale is performed to characterise the interaction between vibrational relaxation and instability growth.


The streamwise evolution of the Damk\"{o}hler number at five characteristic stations is presented in \ref{figure:damkohlernumber}. The Damk\"{o}hler number, defined as the ratio of the viscous timescale $\delta^{*}/u_{\tau}$ to the vibrational relaxation timescale, quantifies the relative importance of instability-driven energy redistribution compared to molecular relaxation. Across all thermal conditions, a consistent trend emerges: the Damk\"{o}hler number peaks are progressively confined towards the wall during the secondary instability region. This wall-confinement indicates that vibrational excitation localises where the viscous timescale approaches the relaxation timescale. The magnitude of these peaks varies systematically with Mach number and wall temperature. At Mach~6 with $T_w = 1200$~K (panel~(a)), moderate peaks emerge, whilst elevated wall temperature (panel~(b)) amplifies them significantly. At Mach~8 (panels~(c) and~(d)), peak magnitudes are substantially reduced due to faster viscous timescales, with cold wall conditions (panel~(c)) showing the most suppression. These results contrast with the laminar findings of Chapter~4, demonstrating that the turbulent regime exhibits different dynamics where the balance between viscous timescale compression and enhanced vibrational energy reshapes the spatial confinement and intensity of thermal nonequilibrium.


Finally, the vibrational temperature deviation profiles presented in \ref{figure:iles_tmtv_profiles} across both linear and semi-logarithmic coordinates reveal the wall-normal progression of nonequilibrium through the streamwise stations. Consistent with the Damköhler number analysis, the vibrational temperature deficit progresses towards the wall as the flow transitions through the secondary instability region. In the laminar station, the vibrational temperature deviation remains distributed throughout the boundary layer, but as secondary instabilities develop, the profiles sharpen and concentrate their maximum gradient in the wall-adjacent region. This progressive wall-confinement reflects the same timescale balance discussed previously, nonequilibrium effects intensify where viscous processes and molecular relaxation become comparable. The semi-logarithmic coordinates (right column) reveal the inner-layer structure more clearly, exposing pronounced features in the buffer layer ($y^+ \sim 5$--$30$) and logarithmic layer ($y^+ \sim 30$--$300$) across the thermal conditions. At Mach~6, particularly at elevated wall temperature, the vibrational temperature deviation exhibits a secondary peak in the buffer layer during transition, indicating localised energy redistribution driven by the intense shear and secondary instability dynamics. At Mach~8, these buffer-layer features are substantially suppressed, consistent with the reduced Damköhler number magnitudes, demonstrating that the faster viscous timescale inhibits vibrational excitation regardless of wall temperature effects. The semi-log visualisation thus confirms that the inner-layer dynamics controlling thermal nonequilibrium are fundamentally reshaped by the competing influences of Mach number and wall temperature on the viscous and relaxation timescales.




% \subsection{Second-order statistics}



\section*{Chapter Summary}

This chapter examined how thermal modelling affects hypersonic boundary-layer transition and how Mach number and wall temperature influence transition behaviour in thermal nonequilibrium flows. The first part compared three thermal models, \gls{CPG}, \gls{TEQ}, and \gls{TNE}, using \gls{DNS} at Mach 6 with hot walls. The second part to examined transition across Mach 6 and 8 with cold and hot walls, isolating the independent effects of Mach number and wall temperature.

The \gls{DNS} results at Mach 6 with $T_w = 1800$ K showed that thermal modelling substantially altered transition characteristics. \gls{TEQ} exhibited the highest second-mode growth rates followed by \gls{CPG}, whilst \gls{TNE} showed the most attenuated amplification due to vibrational relaxation. \gls{TEQ} transitioned earliest with pronounced overshoot in skin friction and Stanton number, \gls{CPG} transitioned downstream at a slightly delayed location, and \gls{TNE} delayed onset by approximately $150\delta^*$ (19\% of the domain) with subdued overshoot. This delay reflected the stabilising effects of vibrational nonequilibrium base flow has on secondary instability growth. Heat flux organisation and velocity streaks developed earlier and more vigorously in \gls{CPG} and \gls{TEQ}, whilst \gls{TNE} showed delayed and attenuated streak formation. In the fully turbulent regime, Reynolds stress components collapsed across all three models, indicating that turbulence structure was largely model-independent once transition was complete. However, thermal statistics showed that vibrational nonequilibrium persisted much further downstream, with mean temperature and pressure fluctuations remaining elevated compared to perfect-gas predictions.

The comparison of \gls{DNS} and \gls{ILES} demonstrated that \gls{ILES} was a cost-effective alternative, reducing computational cost by 6.4 times whilst preserving transition onset location and early transition dynamics. Boundary-layer integral parameters tracked closely between \gls{DNS} and \gls{ILES}, with \gls{TNE} showing consistent behaviour throughout the domain. Minor differences in the fully turbulent regime reflected underresolution on coarser grids, but overall agreement was sufficient to establish \gls{ILES} as a practical framework for parametric studies. The parametric study across Mach 6 and 8 with cold and hot walls revealed that the Mach 6 cold wall case exhibited the highest second-mode amplification, followed by Mach 6 hot wall, with both Mach 8 cases substantially stabilised. The Mach 6 cold wall case transitioned earliest with pronounced resonance peaks, whilst the Mach 8 hot wall case showed the most extended transition. Thermal nonequilibrium analysis demonstrated that wall-confinement of vibrational excitation was most pronounced at Mach 6 elevated wall temperature, where viscous timescales were slowest relative to relaxation processes. At Mach 8, faster viscous timescales substantially suppressed vibrational excitation throughout the boundary layer regardless of wall temperature. Vibrational temperature deviation profiles confirmed this systematic variation, with semi-logarithmic coordinates revealing enhanced buffer-layer nonequilibrium exclusively at Mach 6 conditions. Streak amplitude evolution showed rapid amplification at Mach 6 with early secondary instability threshold crossings, whilst Mach 8 cases maintained suppressed amplitudes reflecting the inhibition of secondary instability at higher Mach number. Heat flux analysis revealed that vibrational heat flux contributions were minimal at Mach 8 cold conditions but substantial at Mach 6 elevated wall temperature, demonstrating the strong temperature dependence of vibrational excitation. These findings showed that the balance between viscous timescale compression at higher Mach numbers and enhanced vibrational energy at elevated wall temperatures fundamentally reshaped transition dynamics and thermal nonequilibrium intensity across hypersonic flight regimes.

% \section*{Chapter Summary}

% This chapter examined the influence of thermal modelling on hypersonic boundary-layer transition and assessed how Mach number and wall temperature modulated transition behaviour in thermal nonequilibrium flows. The first part of the chapter focused on comparing three thermal models---\gls{CPG}, \gls{TEQ}, and \gls{TNE}---through \gls{DNS} at Mach 6 with elevated wall temperature, before a grid refinement study validated the use of \gls{ILES} for parametric investigations. The second part employed \gls{ILES} to systematically examine transition across four conditions spanning Mach 6 and 8 with cold and hot walls, isolating the independent effects of Mach number and wall temperature on thermal nonequilibrium dynamics.

% \gls{DNS} of boundary-layer transition at Mach 6 with hot wall ($T_w = 1800$ K) revealed that thermal modelling substantially altered transition characteristics despite preserving underlying instability mechanisms. Linear stability analysis predicted that \gls{TEQ} exhibited the highest second-mode growth rates, followed by \gls{CPG}, with \gls{TNE} showing the most attenuated amplification due to vibrational relaxation extracting energy from acoustic modes. Nonlinear simulations confirmed this hierarchy: \gls{TEQ} transitioned earliest with pronounced overshoot in skin friction and Stanton number, \gls{CPG} transitioned downstream at a slightly delayed location, whilst \gls{TNE} displayed delayed onset by approximately $150\delta^*$ (19\% of the domain) with considerably subdued overshoot. This delay reflected stabilising effects of vibrational nonequilibrium on secondary instability growth. Spanwise heat-flux organisation and coherent velocity streaks developed earlier and more vigorously in \gls{CPG} and \gls{TEQ}, whilst \gls{TNE} exhibited delayed and attenuated streak formation, indicating suppression of secondary instabilities through vibrational energy sequestration. Wall-normal disturbance kinetic energy distributions revealed that \gls{TNE} concentrations remained confined near the wall initially before extending into the boundary layer at downstream locations, reflecting the delayed amplification characteristic of thermally nonequilibrium flows. In the fully turbulent regime, Reynolds stress components collapsed across all three models, demonstrating that turbulence structure itself was largely model-independent once transition was complete. However, thermal statistics revealed that vibrational nonequilibrium persisted much further downstream, with mean temperature and pressure--temperature fluctuations remaining elevated compared to perfect-gas predictions.

% A grid study comparing fine-resolution \gls{DNS} ($\Delta x^+ \approx 10.66$, $\Delta z^+ \approx 4.95$) against coarser \gls{ILES} ($\Delta x^+ \approx 23$--$26$, $\Delta z^+ \approx 14$--$17$) grids demonstrated the feasibility of implicit large eddy simulation as a cost-effective alternative to \gls{DNS} for transition studies. \gls{ILES} simulations reduced computational cost by approximately 6.4 times whilst preserving transition onset location and early transition dynamics across all three thermal models. Boundary-layer integral parameters including displacement and momentum thicknesses tracked closely between \gls{DNS} and \gls{ILES} across all cases, with \gls{TNE} showing remarkably consistent behaviour throughout the domain. Minor differences in the fully turbulent regime emerged in buffer-layer velocity profiles and turbulent Mach number, reflecting underresolution of fine-scale coherent structures on coarser grids. However, the overall agreement was sufficient to establish \gls{ILES} as a practical framework for systematic parametric studies where \gls{DNS} resolution was computationally prohibitive.

% The combined influence of Mach number and wall temperature on transition in thermal nonequilibrium flows was assessed through \gls{ILES} across four conditions: Mach 6 and 8 with cold ($T_w = 1200$ K) and hot ($T_w = 1800$ K) walls. Linear stability analysis confirmed that the Mach 6 cold wall case exhibited the highest second-mode amplification ($N\text{-factor} \approx 8$), followed by Mach 6 hot wall ($N \approx 6$), with both Mach 8 cases substantially stabilised ($N \approx 4$ and $2$). This progression manifested in transition characteristics: the Mach 6 cold wall case transitioned earliest with pronounced resonance peaks in skin friction and Stanton number, whilst the Mach 8 hot wall case exhibited the most extended transition process with subdued linear growth peaks. Damköhler number analysis demonstrated that wall-confinement of vibrational excitation was most pronounced at Mach 6 elevated wall temperature, where viscous timescales were slowest relative to relaxation. At Mach 8, faster viscous timescales substantially reduced Damköhler number magnitudes regardless of wall temperature, suppressing vibrational excitation throughout. Vibrational temperature deviation profiles confirmed this systematic variation, with semi-logarithmic coordinates revealing enhanced buffer-layer nonequilibrium exclusively at Mach 6 conditions. Streak amplitude evolution showed rapid amplification at Mach 6 with secondary instability thresholds crossed early, particularly in cold wall conditions, whilst Mach 8 cases maintained suppressed streak amplitudes reflecting weaker secondary instability growth at higher compressibility. Heat flux decomposition revealed that vibrational heat flux contributions were minimal at Mach 8 cold conditions but substantial at Mach 6 elevated wall temperature, demonstrating pronounced temperature dependence of vibrational excitation strength. These findings established that the balance between viscous timescale compression at higher Mach numbers and enhanced vibrational energy at elevated wall temperatures fundamentally reshaped transition dynamics and thermal nonequilibrium intensity across hypersonic flight regimes, with significant implications for aerothermal design of high-speed vehicles.
